% Pregunta 877 (2003-julio-OEP2003, n.º 5): funciones de reacción en cantidades (decrecientes: sustitutos estratégicos) y en precios (crecientes: complementos estratégicos). Autor: Víctor Gutiérrez Marcos
\documentclass[tikz,border=6pt]{standalone}
\input{tikz-tcee}
\usetikzlibrary{calc}
\begin{document}
\begin{tikzpicture}[tcee,x=1.1cm,y=1.1cm]
  % ---------- Panel 1: cantidades (Cournot) ----------
  % R1: q1 = (4 - 0,8 q2)/2 ; R2: q2 = (4 - 0,8 q1)/2 ; equilibrio q = 4/2,8 = 1,43
  \node[font=\small\bfseries] at (2.5,5.3) {Cantidades (Cournot)};
  \ejes{$q_1$}{$q_2$}{5.2}{4.9}
  \draw[curva1] (0.16,4.6) -- (2,0) node[etiqueta,below right] {$R_1$};
  \draw[curva2] (0,2) node[etiqueta,above right] {$R_2$} -- (4.6,0.16);
  \proyeccion{(1.43,1.43)}{$q^C$}{$q^C$}
  \node[punto] at (1.43,1.43) {};
  \draw[flecha] (1.56,1.91) -- (1.08,3.11);
  \node[anchor=west,align=left,font=\small] at (1.35,3.3) {si $q_2$ sube,\\ $q_1$ baja};
  \node[align=center,font=\small] at (2.5,-1.05) {$R_1$ decreciente:\\ \textbf{sustitutos estratégicos} (b)};

  % ---------- Panel 2: precios (Bertrand) ----------
  % R1: p1 = 1,2 + 0,4 p2 ; R2: p2 = 1,2 + 0,4 p1 ; equilibrio p = 1,2/0,6 = 2
  \begin{scope}[xshift=7.6cm]
  \node[font=\small\bfseries] at (2.5,5.3) {Precios (Bertrand)};
  \ejes{$p_1$}{$p_2$}{5.2}{4.9}
  \draw[curva1] (1.2,0) -- (3.04,4.6) node[etiqueta,right] {$R_1$};
  \draw[curva2] (0,1.2) -- (4.9,3.16) node[etiqueta,right] {$R_2$};
  \proyeccion{(2,2)}{$p^B$}{$p^B$}
  \node[punto] at (2,2) {};
  \draw[flecha] (1.92,2.61) -- (2.52,4.11);
  \node[anchor=east,align=right,font=\small] at (2.05,3.7) {si $p_2$ sube,\\ $p_1$ sube};
  \node[align=center,font=\small] at (2.5,-1.05) {$R_1$ creciente:\\ \textbf{complementos estratégicos}};
  \end{scope}
\end{tikzpicture}
\end{document}
